\section{Numerical Verification}\label{app:numerics}

We checked every result in \cref{app:proofs} against solutions computed by brute force from the primitives of the model. The scripts are in \texttt{code/checks}, run with \texttt{python3 -I} and finish in under thirty seconds each. No solver uses the formula it checks: equilibria come from enumeration, backward induction or direct optimization, and closed forms are computed separately. Where a statement depends on a tie-breaking rule, the script solves the game under that rule or computes every equilibrium outcome. \texttt{code/checks/README.md} maps the names used in the script comments to the labels used here.

\paragraph{Bidding.} \texttt{v2\_m1\_lemma1.py} works in exact rational arithmetic. In the clock auction (300 lattice draws) dropping at one's maximal credible bid is weakly dominant in every draw, and Europe's payoff equals $\pos{\min\{Q,V,\bar m_U\}}$ exactly. In the alternating-raise auction, 800 games are solved for the full set of pure subgame-perfect outcomes with no tie-breaking. The winner matches \cref{lem:bidding} in all 580 games with $Q$ more than one grid step below $\bar z_U$ and in all 44 with $Q$ more than one step above it. The largest payoff to Europe is within one grid step of the formula in all 800, and in 108 games the smallest is lower by more than a step, the price indeterminacy of the claim in \cref{app:proofs}. When the losing bidder raises whenever indifferent, the price is within one grid step of $\pos{Q}$.

\paragraph{Bargaining.} \texttt{v2\_m1\_prop1.py} runs finite-horizon backward induction in which each proposer compares every offer on a 601-point grid with the payoff of an unacceptable offer. Over 500 runs the largest gap between Europe's value and $x^*(\Delta)$ is $1.7\times10^{-11}$, and with $\lambda>0$ no proposer weakly prefers delay. Richardson extrapolation from $\Delta=0.004,0.002,0.001$ on 120 draws gives a largest error against $P^*$ of $7.5\times10^{-6}$. Iterating the supremum and infimum maps of Step 4 from $[-10V,10V]^2$ converges to $(x^*,x^*,-x^*,-x^*)$ within $4.4\times10^{-12}$ (300 draws). Parts (a) to (c) of \cref{cor:master} hold in every draw, finite differences in $V$ match $\beta\psi$ to $6.6\times10^{-10}$, and $\hat Q$ classifies every draw. For the timing of the footnote to \cref{prop:master}, the stationary value matches its formula to $1.6\times10^{-14}$ (300 draws), and the nonstationary profile that gives Europe close to $V$ passes a one-shot deviation check in all 67 draws in which it was built.

\paragraph{Slack principle, rationing and hedges.} \texttt{v2\_m3\_prop2.py} draws one to four instruments and one to four demands, with additive or general monotone set functions, premia of every sign, reversible and sunk demands, zero or positive execution costs, and integer draws that create ties. The full set of subgame-perfect outcomes of the adaptive protocol comes from set-valued backward induction over states, and a separate solver finds the best single ultimatum. $B(J)$ and $\min\{T,\Lambda\}$ match to $8.9\times10^{-16}$ over 3,000 instances, and the acceptance rule of \cref{prop:slack}(b) has no mismatch in 123,300 pairs. The bound of part (c) is never violated in 4,000 adaptive instances, and Washington's payoff never exceeds $V_1$ (it falls strictly below in 325 of 2,000). Where a sufficient condition of part (d) holds (1,593 instances), the adaptive and single-ultimatum values coincide, and the ratchet count $n\bmod(b+1)$ matches in all 300 draws. In the burning example (two instruments with $u=1.2$ and $1.0$ and $d=0$, $Q=-5$, sunk demands with $t=1$, $w=5$ and $t=1.5$, $w=1$, execution cost $0.1$ each), Washington first issues an ultimatum that Europe refuses in every equilibrium, which burns one instrument, and then obtains the first demand, for a payoff of $4.9$ against $5$ under a single ultimatum. When a refusal can make a new instrument credible, Europe's loss exceeds the bite of the initial credible set by up to $4.91$ but never that of the enlarged set, which is why the proposition fixes $J^c$. A plain knapsack with capacity $\min\{\Lambda,B(J^c)\}$ misses Washington's optimum in 521 of 1,046 reversible-only instances, and the program of \cref{cor:rationing} in none. The linkage condition of \cref{prop:counter}(b) has no mismatch in 3,000 draws.

\paragraph{Rupture risk.} \texttt{v2\_rupture\_bite.py} computes the bite of a withdrawal directly from $P^*$ of \cref{prop:master} in 13,079 random draws with $Q+d<V$ ($\lambda=0$ in about half). The bound of \cref{rem:rupture} has no violation, at $\lambda=0$ the bite agrees with \cref{prop:slack}(a) to $8.9\times10^{-16}$, and the example of the text gives $\Lambda=1$, $\Lambda_\psi=2.2$ and a bite of $1.8$.

\paragraph{Counter-threats.} \texttt{v2\_m3\_prop2e.py} lets China choose counter-threats from a grid before Washington's single ultimatum. Over 4,000 instances the acceptance rule of \cref{prop:counter}(a)(i) has no mismatch in 13,286 checks. China gains in 615 instances, each time by changing Washington's demand set; a counter-threat is executed on the path in 140 instances, and each executed threat crowds out another demand. Washington never gains. Europe's loss rises in 89 instances and falls in 405, and with a divisible sunk demand that China cannot threaten it equals $B(J^c)$ to $8.9\times10^{-16}$.

\paragraph{Conditionality.} \texttt{v2\_m3\_prop3.py} draws 1,500 problems with $H$ from six convex families and maximizes each on a grid of 200,001 points with golden-section polishing. The formula for $\kappa^*$ is never beaten (largest shortfall $8.9\times10^{-16}$), all five regimes occur, and finite differences in $Q_0$ give the derivatives of \cref{prop:conditionality}(c). For part (d), 3,000 draws per distribution show $\bar\kappa$ rising with $s_X$ in 1,218 normal, 1,198 logistic, 974 exponential and 973 uniform cases, and the derivative matches finite differences to $3\times10^{-11}$.

\paragraph{Preventive control.} \texttt{v2\_m4\_prop4.py} solves the two-period game by backward induction from primitives on nine edge cases, 20,000 integer draws with dyadic $\delta$ (so ties are exact) and 10,000 continuous draws. The issue decision never differs from \cref{prop:control}; the errors in $\Gamma$, $T^*$ and Europe's payoff are zero on integer draws and below the grid mesh otherwise; and none of the 7,715 draws with an ultimatum has a refusal on the path. The identity of \cref{cor:control}(b) holds to $2.7\times10^{-15}$, the purchase without bite gives $\Gamma=1$ and $T^*=2.5$, $\Gamma$ never falls as $\Delta_c$ rises (3,000 draws), and the special case after \cref{prop:control} matches in 1,504 draws. \texttt{v2\_m4\_prop4\_china.py} adds China's announcement of $r^c$ (step $\tfrac12$) and a stake, with $e^C(r^c)=\gamma r^c$. In 5,000 integer draws every China-optimal announcement satisfies \cref{cor:control}(d)(i) to (iii), $z_C=0$ leaves the outcome open in 183 draws, as (iv) allows, and a self-rewarding punishment ($\gamma=-\tfrac12$) is carried out on the path in 968.

\paragraph{Cuts and windows.} \texttt{v2\_m4\_cuts\_windows.py}, part A, draws 20,000 games with one to five dates. In the 13,927 draws inside the scope of \cref{lem:selflimit} the largest errors are $3.6\times10^{-15}$ for the bite and $1.8\times10^{-15}$ for the cost; in all 6,073 outside it the true loss exceeds the formula. Part B solves 60,000 multi-date games by backward induction in exact rational arithmetic. The condition with $t'>t$ misclassifies 1,030 draws, all with some $u_t>0$, and the condition of \cref{cor:windows} matches every draw. Of the 20,199 games with an ultimatum, 13,923 issue it at the last eligible date and 1,202 at a date that is not itself eligible. No refused ultimatum is ever optimal.

\paragraph{Self-binding.} \texttt{v2\_m4\_selfbind.py} solves 6,000 integer draws with $Q_1<0$, $\Delta_c=\infty$ and no investment. Parts (a), (b) and (d) of \cref{rem:bind} have no mismatch, partial binding invites an ultimatum that would not otherwise be issued in 210 draws, and the example of part (c) is reproduced.

\paragraph{Breakdown.} \texttt{v2\_m4\_prop5.py} maximizes Washington's gain directly over the payment, without the first-order condition, for both kinds of offer, with $\sigma$ drawn from six families (1,500 draws each): uniform, truncated normal, Beta(2,2), triangular, truncated exponential, and a bimodal mixture that is not log-concave. For the five log-concave families the issuance condition of \cref{prop:breakdown}(c) matches every draw away from the threshold, and the largest errors are $3.4\times10^{-7}$ in the cutoff and $2.2\times10^{-7}$ in $P_b$. For the comparative statics (500 draws per family) part (d) has no violation in any family, the bimodal one included; with the purchase option, violations occur only in the directions part (e) says can fail, and the four examples in the proof are reproduced.

\paragraph{Who decides.} \texttt{v2\_m2\_remarks\_GFE.py} checks \cref{rem:gatekeeper} by brute force on a grid of mesh $10^{-3}$ (20,000 draws; largest errors $4.4\times10^{-16}$ in the payment and $8.9\times10^{-16}$ in the slack) and \cref{rem:firm} on 200,000 integer draws, with no mismatch. For \cref{rem:council} it reads the population shares of Council Decision 2025/2427 \citep{council2025population} from \texttt{pop\_check.py}. The first-stage blocking rule matches the treaty rule on all 403,168 sets of members tested and the delegated-act rule on 241,708, and the coalitions of \cref{tab:council} and $\Phi(0)$ are computed exactly. \texttt{council\_revealed\_blocking.py} applies the same rule in integer arithmetic to 2026 populations and the EUROPOP2023 projections for 2029 to 2031. The seven states that \emph{Europe 2031} names hold 28.06 percent of the population in 2026 and 27.76 percent in 2031. In every year Germany, France and Spain are the only states that complete a block alone, and without them and Romania the smallest completions have three states (for example Belgium, Czechia and Portugal). \Cref{sec:council} and \cref{app:scenario} use the population shares and completions. Exact Shapley--Shubik and normalized Banzhaf indices, checked against 200,000 random orderings, give Germany 17.54 and 12.03 percent against 3.66 and 3.80 percent for the Netherlands, ratios of 4.8 and 3.2.

\paragraph{Hold-up.} \texttt{v2\_m2\_prop6\_holdup.py} solves the discrete game by backward induction on the grid of offer dates, with random piecewise-constant cost paths, for 6,000 parameter sets. Where $Q\le\bar Q$, the largest error against $x^\circ(0)$ is $3.1\times10^{-7}$, and the test $Q\le\bar Q$ classifies every draw. Where exit binds, the value exceeds $x^\circ(0)$ in every draw and converges to the stopping value at first order, and the counterexample in the proof is reproduced ($-0.3679$ against $-0.7358$ for $x^\circ(0)$). \texttt{v2\_m2\_interruption.py} solves the stock-buffer game of \cref{cor:interrupt} and the three timing variants by backward induction over the stock (3,000 parameter sets, exit binding without interruption in 977). The largest error against part (ii) is $8.1\times10^{-4}$, the decision agrees with $(1-\psi)D>\pos{Q-x^\circ_0}$ in all 2,834 draws not near a tie, and every draw that interrupts at the full stock interrupts at every lower stock. No draw of variant (a) interrupts before $t_U$; in variant (b) the gain is exactly zero in all 2,852 draws with $\ell>0$; and in variant (c) the gain is positive in 551 draws in which the condition of part (i) fails.
